\documentclass[12pt,a4paper]{article}

\usepackage{fontspec}
\usepackage{amsmath,amssymb,mathtools}
\usepackage{booktabs}
\usepackage{enumitem}
\usepackage{hyperref}

\setmainfont{Times New Roman}
\setlength{\oddsidemargin}{0cm}
\setlength{\evensidemargin}{0cm}
\setlength{\textwidth}{16cm}
\setlength{\topmargin}{-1cm}
\setlength{\textheight}{24.5cm}
\hypersetup{colorlinks=true,linkcolor=blue,urlcolor=blue}

\title{Mô hình MILP cho bài toán Truck--Drone Resupply PDP}
\author{Mô hình một điểm hẹn cho mỗi chuyến drone, xe tải nhiều chuyến}
\date{}

\begin{document}
\maketitle

\section{Mô tả bài toán}

Bài toán xét một đội xe tải và drone xuất phát từ depot \(0\) để phục vụ tập khách
hàng \(N\). Mỗi khách hàng được xe tải ghé đúng một lần. Với khách hàng giao hàng
trực tiếp \(i\in N^D\), kiện hàng của \(i\) có thể được chất lên xe tải tại depot
đầu một chuyến truck, hoặc được drone mang từ depot tới một điểm hẹn trên hành trình
truck. Kiện C1 phải được giao trước khi truck quay lại depot kết thúc chuyến đã nhận
kiện đó. Với cặp pickup--delivery \((p,\ell)\in\mathcal A\), pickup \(p\) và delivery
\(\ell\) phải nằm trên cùng một truck, và delivery phải nằm sau pickup trong toàn
hành trình của truck đó.

Mỗi truck có thể đi nhiều chuyến. Mỗi chuyến truck có dạng
\[
  0 \rightarrow \cdots \rightarrow 0.
\]
Khi truck quay lại depot để bắt đầu chuyến tiếp theo, truck không bắt buộc phải rỗng:
hàng đã pickup ở một khách hàng có thể được giữ trên truck để giao cho delivery tương
ứng ở chuyến sau, miễn là capacity truck luôn được tôn trọng. Lần xuất phát đầu tiên
từ depot không tính thời gian bốc dỡ truck; từ chuyến thứ hai trở đi có tính thời gian
bốc dỡ tại depot. Lần cuối truck quay về depot để kết thúc hành trình không tính thêm
thời gian bốc dỡ.

Mỗi chuyến drone chỉ có một điểm hẹn với một truck tại một khách hàng. Drone xuất
phát từ depot, bay tới điểm hẹn, bàn giao hàng cho truck, rồi quay về depot:
\[
  0 \rightarrow h \rightarrow 0.
\]
Mô hình này không dùng multi-rendezvous cho drone.

\paragraph{Khoảng cách và vận tốc.}
Truck dùng khoảng cách Manhattan, vận tốc \(30\) km/h. Drone dùng khoảng cách
Euclidean, vận tốc \(60\) km/h. Thời gian di chuyển tính theo phút:
\[
  t^T_{ij}=\frac{|x_i-x_j|+|y_i-y_j|}{30}\cdot 60,\qquad
  t^D_{ij}=\frac{\sqrt{(x_i-x_j)^2+(y_i-y_j)^2}}{60}\cdot 60.
\]

\section{Tập chỉ số và tham số}

\begin{table}[h]
\centering
\begin{tabular}{ll}
\toprule
Ký hiệu & Ý nghĩa\\
\midrule
\(0\) & depot\\
\(N\) & tập khách hàng\\
\(N^D\) & tập khách hàng giao hàng trực tiếp, loại \texttt{D}\\
\(N^P\) & tập node pickup, loại \texttt{P}\\
\(N^{DL}\) & tập node delivery của pickup--delivery, loại \texttt{DL}\\
\(\mathcal A\) & tập cặp pickup--delivery \((p,\ell)\)\\
\(K\) & tập truck\\
\(R\) & tập drone\\
\(H=\{1,\ldots,H_{\max}\}\) & tập chuyến tiềm năng của mỗi truck\\
\(P=\{1,\ldots,P_{\max}\}\) & tập vị trí tiềm năng trong một chuyến truck\\
\(S=\{1,\ldots,S_{\max}\}\) & tập chuyến bay tiềm năng của mỗi drone\\
\bottomrule
\end{tabular}
\end{table}

\begin{table}[h]
\centering
\begin{tabular}{ll}
\toprule
Ký hiệu & Ý nghĩa\\
\midrule
\(q_i\) & khối lượng kiện hàng của khách \(i\)\\
\(r_i\) & release date của kiện hàng \(i\)\\
\(Q_T\) & tải trọng truck\\
\(Q_D\) & tải trọng drone\\
\(s_T\) & thời gian truck phục vụ một khách hàng\\
\(\delta_m\) & thời gian bàn giao drone--truck tại điểm hẹn\\
\(\delta_d\) & thời gian drone chất hàng tại depot trước mỗi chuyến bay\\
\(\delta_T\) & thời gian truck bốc dỡ tại depot trước chuyến truck thứ hai trở đi\\
\(E\) & endurance tối đa của một chuyến bay drone\\
\(M\) & hằng số Big-M đủ lớn\\
\(t^T_{ij}\) & thời gian truck đi từ \(i\) đến \(j\)\\
\(t^D_{ij}\) & thời gian drone bay từ \(i\) đến \(j\)\\
\bottomrule
\end{tabular}
\end{table}

\section{Biến quyết định}

\subsection{Route truck}

\[
y_{kmpi}=
\begin{cases}
1, & \text{nếu truck }k\text{ phục vụ khách }i\text{ ở vị trí }p
\text{ trong chuyến }m,\\
0, & \text{ngược lại,}
\end{cases}
\]
\[
u_{kmp}=
\begin{cases}
1, & \text{nếu vị trí }p\text{ của chuyến }m\text{ của truck }k\text{ được dùng,}\\
0, & \text{ngược lại,}
\end{cases}
\]
\[
g_{km}=
\begin{cases}
1, & \text{nếu chuyến }m\text{ của truck }k\text{ được dùng,}\\
0, & \text{ngược lại.}
\end{cases}
\]
\[
z_{kmpij}=
\begin{cases}
1, & \text{nếu trong chuyến }m\text{ của truck }k,\ i\text{ ở vị trí }p
\text{ và }j\text{ ở }p+1,\\
0, & \text{ngược lại,}
\end{cases}
\]
\[
\lambda_{kmpi}=
\begin{cases}
1, & \text{nếu }i\text{ là khách cuối của chuyến }m\text{ của truck }k
\text{ tại vị trí }p,\\
0, & \text{ngược lại.}
\end{cases}
\]

\subsection{Nguồn hàng và drone}

\[
a_{kmi}=
\begin{cases}
1, & \text{nếu kiện C1 }i\text{ được chất lên truck }k\text{ tại depot đầu chuyến }m,\\
0, & \text{ngược lại,}
\end{cases}
\]
\[
b_{kmpi}=
\begin{cases}
1, & \text{nếu kiện C1 }i\text{ được drone giao cho truck }k
\text{ tại vị trí }p\text{ của chuyến }m,\\
0, & \text{ngược lại,}
\end{cases}
\]
\[
o_{kmpi}=
\begin{cases}
1, & \text{nếu node tại vị trí }p\text{ là }i\text{ và kiện }i
     \text{ được nhận từ drone tại chính vị trí đó},\\
0, & \text{ngược lại.}
\end{cases}
\]
Biến phụ $o_{kmpi}$ được tuyến tính hóa bởi
\begin{equation}
o_{kmpi}\le y_{kmpi},\qquad
o_{kmpi}\le b_{kmpi},\qquad
o_{kmpi}\ge y_{kmpi}+b_{kmpi}-1.
\end{equation}
\[
\rho_{kmp}=
\begin{cases}
1, & \text{nếu có một chuyến drone hẹn truck }k\text{ tại vị trí }p
\text{ của chuyến }m,\\
0, & \text{ngược lại.}
\end{cases}
\]
\[
w_{dskmpj}=
\begin{cases}
1, & \text{nếu chuyến bay }s\text{ của drone }d\text{ gặp truck }k
      \text{ trong chuyến }m,\\
   & \text{tại vị trí }p\text{, node }j,\\
0, & \text{ngược lại,}
\end{cases}
\]
\[
\eta_{ds}=
\begin{cases}
1, & \text{nếu chuyến bay }s\text{ của drone }d\text{ được dùng,}\\
0, & \text{ngược lại.}
\end{cases}
\]

\subsection{Thời gian và tải}

\[
T^0_{km}\ge0,\quad A_{kmp}\ge0,\quad D_{kmp}\ge0,\quad B_{km}\ge0,
\]
trong đó \(T^0_{km}\) là thời điểm truck \(k\) rời depot bắt đầu chuyến \(m\),
\(A_{kmp}\), \(D_{kmp}\) là thời điểm truck đến/rời vị trí \(p\), và \(B_{km}\)
là thời điểm truck quay về depot kết thúc chuyến \(m\).

\[
F^{dep}_{ds}\ge0,\quad F^{arr}_{ds}\ge0,\quad F^{ret}_{ds}\ge0,\quad C_{\max}\ge0.
\]
\[
L_{kmp}\ge0,\quad L^{end}_{km}\ge0
\]
là tải trên truck \(k\) ngay sau khi rời vị trí \(p\) trong chuyến \(m\). Ký hiệu
\(L_{km0}\) là tải truck ngay khi rời depot đầu chuyến \(m\), và \(L^{end}_{km}\)
là tải còn trên truck khi chuyến \(m\) quay về depot. Tải còn lại này chỉ có thể là
hàng pickup C2 chưa tới delivery tương ứng.

\section{Hàm mục tiêu}

\[
  \min C_{\max}.
\]

\section{Ràng buộc}

\subsection{Cấu trúc route truck nhiều chuyến}

\begin{equation}
  \sum_{i\in N} y_{kmpi}=u_{kmp},
  \qquad \forall k\in K,\ m\in H,\ p\in P.
\end{equation}

\begin{equation}
  u_{kmp}\le g_{km},
  \qquad \forall k\in K,\ m\in H,\ p\in P.
\end{equation}

\begin{equation}
  u_{km,p+1}\le u_{kmp},
  \qquad \forall k\in K,\ m\in H,\ p=1,\ldots,P_{\max}-1.
\end{equation}

\begin{equation}
  g_{k,m+1}\le g_{km},
  \qquad \forall k\in K,\ m=1,\ldots,H_{\max}-1.
\end{equation}

\begin{equation}
  \sum_{k\in K}\sum_{m\in H}\sum_{p\in P} y_{kmpi}=1,
  \qquad \forall i\in N.
\end{equation}

\begin{equation}
  z_{kmpij}\le y_{kmpi},\quad
  z_{kmpij}\le y_{km,p+1,j},\quad
  z_{kmpij}\ge y_{kmpi}+y_{km,p+1,j}-1,
\end{equation}
\[
  \forall k\in K,\ m\in H,\ p=1,\ldots,P_{\max}-1,\ i,j\in N.
\]

\begin{equation}
  \lambda_{kmpi}\le y_{kmpi},\quad
  \lambda_{kmpi}\le u_{kmp}-u_{km,p+1},\quad
  \lambda_{kmpi}\ge y_{kmpi}+u_{kmp}-u_{km,p+1}-1,
\end{equation}
\[
  \forall k\in K,\ m\in H,\ p=1,\ldots,P_{\max}-1,\ i\in N.
\]
\begin{equation}
  \lambda_{km,P_{\max},i}=y_{km,P_{\max},i},
  \qquad \forall k\in K,\ m\in H,\ i\in N.
\end{equation}

\subsection{Pickup--delivery}

Pickup và delivery của cùng một cặp phải nằm trên cùng một truck, nhưng có thể nằm
ở hai chuyến khác nhau:
\begin{equation}
  \sum_{m\in H}\sum_{p\in P}y_{kmp p'}
  =
  \sum_{m\in H}\sum_{p\in P}y_{kmp \ell},
  \qquad \forall k\in K,\ (p',\ell)\in\mathcal A.
\end{equation}

Pickup phải đứng trước delivery trong toàn hành trình của truck. Thứ tự toàn hành
trình của vị trí \((m,p)\) được tính bởi \((m-1)P_{\max}+p\):
\begin{equation}
\begin{aligned}
  &\sum_{m\in H}\sum_{p\in P}\big((m-1)P_{\max}+p\big)y_{kmp p'}+1\\
  &\le
  \sum_{m\in H}\sum_{p\in P}\big((m-1)P_{\max}+p\big)y_{kmp\ell}\\
  &\quad +M\left(1-\sum_{m\in H}\sum_{p\in P}y_{kmp p'}\right),
\end{aligned}
\end{equation}
\[
  \forall k\in K,\ (p',\ell)\in\mathcal A.
\]

\subsection{Nguồn hàng cho khách C1}

\begin{equation}
  a_{kmi}+\sum_{p\in P}b_{kmpi}
  =
  \sum_{p\in P}y_{kmpi},
  \qquad \forall k\in K,\ m\in H,\ i\in N^D.
\end{equation}

Do cả nguồn hàng và lần phục vụ đều mang cùng chỉ số chuyến \(m\), phương trình trên
buộc mọi kiện C1 được chất tại depot hoặc nhận từ drone phải được giao trong chính
chuyến truck đó; kiện C1 không thể được mang qua một lần quay về depot.

\begin{equation}
  b_{kmpi}\le \sum_{\substack{q\in P\\q\ge p}} y_{kmqi},
  \qquad \forall k\in K,\ m\in H,\ p\in P,\ i\in N^D.
\end{equation}

\begin{equation}
  b_{kmpi}\le \rho_{kmp},
  \qquad \forall k\in K,\ m\in H,\ p\in P,\ i\in N^D.
\end{equation}

\begin{equation}
  \rho_{kmp}\le u_{kmp},
  \qquad \forall k\in K,\ m\in H,\ p\in P.
\end{equation}

\subsection{Drone capacity và gán chuyến drone}

\begin{equation}
  \sum_{i\in N^D}q_i b_{kmpi}\le Q_D\rho_{kmp},
  \qquad \forall k\in K,\ m\in H,\ p\in P.
\end{equation}

\begin{equation}
  \sum_{d\in R}\sum_{s\in S}\sum_{j\in N} w_{dskmpj}
  =
  \rho_{kmp},
  \qquad \forall k\in K,\ m\in H,\ p\in P.
\end{equation}
\begin{equation}
  w_{dskmpj}\le y_{kmpj},
  \qquad \forall d\in R,\ s\in S,\ k\in K,\ m\in H,\ p\in P,\ j\in N.
\end{equation}

\begin{equation}
  \sum_{k\in K}\sum_{m\in H}\sum_{p\in P}\sum_{j\in N}w_{dskmpj}
  =\eta_{ds},
  \qquad \forall d\in R,\ s\in S.
\end{equation}

\begin{equation}
  \eta_{d,s+1}\le \eta_{ds},
  \qquad \forall d\in R,\ s=1,\ldots,S_{\max}-1.
\end{equation}

\subsection{Truck capacity và tải chuyển tiếp giữa các chuyến}

\begin{equation}
  L_{k10}=\sum_{i\in N^D}q_i a_{k1i},
  \qquad \forall k\in K.
\end{equation}

Từ chuyến thứ hai trở đi, tải đầu chuyến gồm hàng pickup C2 còn trên truck từ chuyến
trước cộng với hàng C1 mới được chất ở depot:
\begin{equation}
  L_{km0}=L^{end}_{k,m-1}+\sum_{i\in N^D}q_i a_{kmi},
  \qquad \forall k\in K,\ m=2,\ldots,H_{\max}.
\end{equation}

\begin{equation}
  L_{km0}\le Q_Tg_{km},
  \qquad \forall k\in K,\ m\in H.
\end{equation}

Tải truck sau khi hoàn tất phục vụ và bàn giao tại vị trí \(p\):
\begin{equation}
  L_{kmp}
  =
  L_{km,p-1}
  +\sum_{i\in N^P}q_i y_{kmpi}
  -\sum_{i\in N^{DL}}q_i y_{kmpi}
  -\sum_{i\in N^D}q_i y_{kmpi}
  +\sum_{i\in N^D}q_i b_{kmpi},
\end{equation}
\[
  \forall k\in K,\ m\in H,\ p\in P.
\]

\begin{equation}
  0\le L_{kmp}\le Q_Tu_{kmp},
  \qquad \forall k\in K,\ m\in H,\ p\in P.
\end{equation}

Tải cuối chuyến bằng tải tại khách cuối cùng của chuyến:
\begin{equation}
  L^{end}_{km}\ge L_{kmp}-M(1-\lambda_{kmpi}),
  \qquad \forall k\in K,\ m\in H,\ p\in P,\ i\in N.
\end{equation}
\begin{equation}
  L^{end}_{km}\le L_{kmp}+M(1-\lambda_{kmpi}),
  \qquad \forall k\in K,\ m\in H,\ p\in P,\ i\in N.
\end{equation}
\begin{equation}
  0\le L^{end}_{km}\le Q_Tg_{km},
  \qquad \forall k\in K,\ m\in H.
\end{equation}

Để cấm tường minh hàng C1 đi qua depot nhưng vẫn cho phép hàng C2 đi qua depot, tải
cuối chuyến phải bằng đúng cân bằng pickup--delivery C2 lũy kế:
\begin{equation}
  L^{end}_{km}
  =
  \sum_{(p',\ell)\in\mathcal A}q_{p'}
  \left(
    \sum_{r=1}^{m}\sum_{s\in P}y_{krs p'}
    -
    \sum_{r=1}^{m}\sum_{s\in P}y_{krs\ell}
  \right),
  \qquad \forall k\in K,\ m\in H.
\end{equation}
Vì pickup và delivery cùng truck và pickup đứng trước delivery, vế phải không âm.
Nếu \(L^{end}_{km}>0\), toàn bộ phần tải đó là hàng C2 đã pickup nhưng chưa tới
delivery tương ứng. Do đó C2 có thể được giữ qua depot, còn C1 thì không.

\subsection{Thời gian truck}

\begin{equation}
  y_{km1i}=1\ \Longrightarrow\
  A_{km1}=T^0_{km}+t^T_{0i},
  \qquad \forall k\in K,\ m\in H,\ i\in N.
\end{equation}

\begin{equation}
  z_{kmpij}=1\ \Longrightarrow\
  A_{km,p+1}=D_{kmp}+t^T_{ij},
\end{equation}
\[
  \forall k\in K,\ m\in H,\ p=1,\ldots,P_{\max}-1,\ i,j\in N.
\]

\begin{equation}
  D_{kmp}\ge A_{kmp}+s_Tu_{kmp}+\delta_m\rho_{kmp},
  \qquad \forall k\in K,\ m\in H,\ p\in P.
\end{equation}
Khi không có drone resupply tại vị trí đang dùng, thời điểm rời được khóa đúng bằng
thời điểm đến cộng service time:
\begin{equation}
  D_{kmp}\le A_{kmp}+s_T+M\rho_{kmp}+M(1-u_{kmp}),
  \qquad \forall k\in K,\ m\in H,\ p\in P.
\end{equation}
\begin{equation}
  D_{kmp}\ge F^{arr}_{ds}+\delta_m
  -M\left(1-\sum_{j\in N}w_{dskmpj}\right),
\end{equation}
\[
  \forall d\in R,\ s\in S,\ k\in K,\ m\in H,\ p\in P.
\]

Thứ tự service--handover phụ thuộc vào nguồn của kiện giao tại node hẹn $j$.
Nếu truck đã có sẵn kiện để phục vụ $j$, service phải hoàn tất trước khi drone tới:
\begin{equation}
 F^{arr}_{ds}\ge A_{kmp}+s_T
 -M\left(1-w_{dskmpj}+o_{kmpj}\right),
\end{equation}
\[
 \forall d\in R,\ s\in S,\ k\in K,\ m\in H,\ p\in P,\ j\in N^D.
\]
Nếu chính kiện $j$ được drone mang tới, truck nhận hàng trước rồi mới service:
\begin{equation}
 D_{kmp}\ge F^{arr}_{ds}+\delta_m+s_T
 -M\left(2-w_{dskmpj}-o_{kmpj}\right),
\end{equation}
\[
 \forall d\in R,\ s\in S,\ k\in K,\ m\in H,\ p\in P,\ j\in N^D.
\]
Đối với node pickup/delivery C2, drone không mang kiện C2 nên luôn thuộc trường hợp
service trước handover.
\begin{equation}
 F^{arr}_{ds}\ge A_{kmp}+s_T-M(1-w_{dskmpj}),
 \quad j\in N^P\cup N^L.
\end{equation}

\begin{equation}
  \lambda_{kmpi}=1\ \Longrightarrow\
  B_{km}=D_{kmp}+t^T_{i0},
  \qquad \forall k\in K,\ m\in H,\ p\in P,\ i\in N.
\end{equation}

Chuyến đầu tiên không cộng thời gian bốc dỡ truck tại depot:
\begin{equation}
  T^0_{k1}\ge r_i a_{k1i},
  \qquad \forall k\in K,\ i\in N^D.
\end{equation}

Từ chuyến thứ hai, truck chỉ rời depot sau khi chuyến trước kết thúc và đã bốc dỡ:
\begin{equation}
  T^0_{km}\ge B_{k,m-1}+\delta_T-M(1-g_{km}),
  \qquad \forall k\in K,\ m=2,\ldots,H_{\max}.
\end{equation}

Release date cho hàng C1 chất lên truck tại depot đầu chuyến:
\begin{equation}
  T^0_{km}\ge r_i a_{kmi},
  \qquad \forall k\in K,\ m\in H,\ i\in N^D.
\end{equation}

\subsection{Release date, drone và đồng bộ}

\begin{equation}
  F^{dep}_{ds}\ge r_i+\delta_d
  -M\left(1-b_{kmpi}\right)
  -M\left(1-\sum_{j\in N}w_{dskmpj}\right),
\end{equation}
\[
  \forall d\in R,\ s\in S,\ k\in K,\ m\in H,\ p\in P,\ i\in N^D.
\]

\begin{equation}
  F^{dep}_{d,s+1}\ge F^{ret}_{ds}+\delta_d-M(1-\eta_{d,s+1}),
  \qquad \forall d\in R,\ s=1,\ldots,S_{\max}-1.
\end{equation}

\begin{equation}
  w_{dskmpj}=1\ \Longrightarrow\
  F^{arr}_{ds}=F^{dep}_{ds}+t^D_{0j},
\end{equation}
\[
  \forall d\in R,\ s\in S,\ k\in K,\ m\in H,\ p\in P,\ j\in N.
\]

\begin{equation}
  w_{dskmpj}=1\ \Longrightarrow\
  F^{ret}_{ds}=F^{arr}_{ds}+\delta_m+t^D_{j0},
\end{equation}
\[
  \forall d\in R,\ s\in S,\ k\in K,\ m\in H,\ p\in P,\ j\in N.
\]

Drone không cần chờ truck tại điểm hẹn. Vì vậy nếu chuyến bay \((d,s)\) được gán
cho điểm hẹn \((k,m,p)\), thời điểm drone tới điểm hẹn không được sớm hơn thời điểm
truck tới vị trí đó:
\begin{equation}
  F^{arr}_{ds}\ge A_{kmp}-M\left(1-\sum_{j\in N}w_{dskmpj}\right),
\end{equation}
\[
  \forall d\in R,\ s\in S,\ k\in K,\ m\in H,\ p\in P.
\]

Truck có thể chờ drone tại điểm hẹn, nên mô hình không dùng tham số giới hạn thời
gian chờ. Để giảm thời gian chờ không cần thiết, nghiệm tối ưu sẽ chọn
\(F^{dep}_{ds}\) phù hợp với release date, thứ tự chuyến bay của drone, và thời điểm
truck tới điểm hẹn; với một lịch cố định, có thể hậu xử lý bằng cách đẩy
\(F^{dep}_{ds}\) muộn nhất có thể mà không làm tăng \(C_{\max}\).

\begin{equation}
  F^{ret}_{ds}-F^{dep}_{ds}\le E+M(1-\eta_{ds}),
  \qquad \forall d\in R,\ s\in S.
\end{equation}

\subsection{Cmax}

\begin{equation}
  C_{\max}\ge B_{km},
  \qquad \forall k\in K,\ m\in H.
\end{equation}

\begin{equation}
  C_{\max}\ge F^{ret}_{ds},
  \qquad \forall d\in R,\ s\in S.
\end{equation}

\subsection{Miền giá trị}

\begin{align*}
y_{kmpi},u_{kmp},g_{km},z_{kmpij},\lambda_{kmpi},
a_{kmi},b_{kmpi},o_{kmpi},\rho_{kmp},w_{dskmpj},\eta_{ds}
&\in\{0,1\},\\
T^0_{km},A_{kmp},D_{kmp},B_{km},
F^{dep}_{ds},F^{arr}_{ds},F^{ret}_{ds},L_{kmp},L^{end}_{km},C_{\max}
&\ge0.
\end{align*}

\section{Ghi chú triển khai}

Mô hình trên và file \texttt{milp.py} cùng sử dụng truck nhiều chuyến, truck và drone
có capacity, C1 phải được nhận và giao trong cùng chuyến truck, hàng pickup C2 có thể
được giữ qua depot để giao ở chuyến sau, và mỗi chuyến drone chỉ có một rendezvous.
Khi CPLEX dừng ở time limit, incumbent có thể chứa slack trong các biến thời gian dù
quyết định route/load/resupply đã cố định. Vì vậy code cố định toàn bộ biến nhị phân
của incumbent và giải lại bài toán thời gian liên tục trước khi báo objective hoặc so
sánh với decoder heuristic.

Khi kiểm tra fitness của một nghiệm heuristic, file encoding lưu thêm drone ID, thứ
tự sortie và các mốc depart/arrive/return. Tất cả quyết định route, depot load,
resupply và gán drone được cố định; phần MILP còn lại chỉ tối ưu các biến thời gian.
Với instance lớn, cùng hệ ràng buộc thời gian được dựng trực tiếp trên các arc và
sortie thực sự xuất hiện trong nghiệm, tránh tạo biến cho các arc không được chọn.
Cách kiểm tra này đã cho cùng $C_{\max}$ trên toàn bộ các bộ 10, 15 và 50 khách
hàng đã thử nghiệm.

\end{document}
